% Generated draft: requires mathematical and source review.
\documentclass[11pt]{article}
\usepackage[margin=25mm]{geometry}
\usepackage{fontspec}
\setmainfont{Latin Modern Roman}
\usepackage{amsmath,amssymb,mathtools,unicode-math}
\setmathfont{Latin Modern Math}
\usepackage{enumitem,needspace,xurl,hyperref}
\setlength{\parindent}{0pt}
\setlength{\parskip}{5pt}
\setlength{\emergencystretch}{3em}
\newcommand{\N}{\mathbb{N}}
\newcommand{\Z}{\mathbb{Z}}
\newcommand{\Q}{\mathbb{Q}}
\newcommand{\R}{\mathbb{R}}
\newcommand{\C}{\mathbb{C}}
\newcommand{\F}{\mathbb{F}}
\newcommand{\A}{\mathbb{A}}
\newcommand{\PP}{\mathbb{P}}
\newcommand{\E}{\mathbb{E}}
\newcommand{\eps}{\varepsilon}
\newcommand{\dd}{\,\mathrm{d}}
\begin{document}
\section*{UnsolvedMath ID: 20003111}
\textbf{Source problem code:} catalogue entry
\subsection{Definitions and mathematical statement}
The source asks whether the Thurston metric can be defined on the non-reduced Teichm\"uller space of a hyperbolic surface with ideal boundary. The admissible surface class, meaning of ideal boundary, boundary marking equivalence, and desired normalization are not specified. Write \(\pi:\mathcal T^{\mathrm{nr}}\to\mathcal T^{\mathrm{red}}\) for the quotient forgetting the additional boundary marking, in the intended reduced/non-reduced conventions.

The aid first reports that ordinary closed-curve lengths and free-endpoint arc lengths factor through \(\pi\). Thus any distance defined solely from those spectra is constant on each boundary-marking fiber and cannot separate non-reduced points in a nontrivial fiber. This is an obstruction for those particular observables, not for arcs with fixed individually marked ideal endpoints.

A second obstruction concerns the raw infimal-Lipschitz formula if its admissible relative maps include all boundary-fixing Lipschitz homeomorphisms. In hyperbolic polar coordinates on the Poincar\'e disk, the radial map
\[
 h_c(r,\theta)=(cr,\theta),\qquad0<c<1,
\]
where \(r\) is hyperbolic distance from the origin, fixes the ideal boundary pointwise and has hyperbolic Lipschitz constant \(c\). Hence the infimum of these constants is zero even for the disk self-comparison, and its unnormalized logarithm collapses to \(-\infty\). These maps are not asserted quasiconformal with a finite global dilatation; restricting the admissible maps to quasiconformal maps would require a separate analysis. The failure pertains to the stated raw formula and class of maps.

For boundary observables, take a hyperbolic \(0\)-metric with boundary defining function \(\rho\), and choose ideal-endpoint geodesics between distinct marked boundary points \(p_i\). A renormalized length has the form
\[
 L_\rho(p_i,p_j)=\lim_{\varepsilon\downarrow0}
 \bigl(\operatorname{length}(\gamma_{ij}\cap\{\rho\geq\varepsilon\})+2\log\varepsilon\bigr),
\]
when that limit and cutoff convention apply. Under \(\rho\mapsto e^\omega\rho\), its change in this convention is \(\omega(p_i)+\omega(p_j)\). The aid's balance criterion for a linear combination is therefore
\[
 \sum_{i<j}c_{ij}L_\rho(p_i,p_j)\text{ is boundary-gauge independent}
 \Longleftrightarrow
 \sum_{j\ne i}c_{ij}=0\quad\text{for every endpoint }i,
\]
with symmetric coefficients \(c_{ij}=c_{ji}\). In the model with at most one chosen length for each distinct unordered endpoint pair, there is no nonzero balanced coefficient system on at most three endpoints over \(\R\); four endpoints admit alternating combinations such as
\[
 L_{12}+L_{34}-L_{14}-L_{23},
\]
which give boundary cross-ratio observables in the appropriate hyperbolic convention. If multiple geodesic classes with the same endpoints are allowed, differences of their lengths already cancel the endpoint gauge, so the four-endpoint minimality statement requires this pair-indexed qualification. The precise \(0\)-metric hypotheses, geodesic choices, and cross-ratio normalization are not supplied. Gauge independence alone is not a distance construction.

Finally, wherever a finite asymmetric arc distance \(d_A\) is defined on the reduced space and the usual non-reduced Teichm\"uller distance \(d_T^{\mathrm{nr}}\) is finite, the aid gives the regularization
\[
 D_\epsilon(X,Y)=d_A(\pi X,\pi Y)+\epsilon d_T^{\mathrm{nr}}(X,Y),\qquad\epsilon>0.
\]
With the usual nonnegativity and triangle inequalities for these inputs, the positive Teichm\"uller term makes this an honest separating asymmetric metric. Its \(\epsilon\)-dependence and use of the preexisting Teichm\"uller metric are retained. The aid describes a canonical pure boundary-sensitive Thurston construction as still open; this regularization does not by itself provide it.
\subsection{Short English statement}
Ordinary length spectra lose the extra boundary marking, and the stated raw Lipschitz formula can collapse if all boundary-fixing radial maps are allowed. Balanced renormalized lengths provide gauge-free boundary data, while adding a positive non-reduced Teichm\"uller term gives a regularized metric. The canonical pure construction remains open according to the aid.
\subsection{Sources}
\begin{itemize}
\item Local UnsolvedMath registry, permanent ID \(20003111\), exact source code \texttt{catalogue entry}; canonical \texttt{aim-topology-notes.json}, \texttt{notes[198]}, tag \texttt{problem}.
\item AIM workshop ``Lipschitz metric on Teichmueller space,'' Problem 3.8: \url{https://aimath.org/pastworkshops/lipschitzteichproblems.pdf}. The supplied verification identifies this official-list question as exact. The PDF is cited without being fetched.
\item The original literature field is empty. No authoritative bibliographic reference for the boundary-marking conventions, hyperbolic \(0\)-metric renormalization, or proposed metric construction is supplied.
\item Research attempt 1: result \texttt{partial\_result}, problem status at run \texttt{open}, rigor label \texttt{proved}, novelty confidence \texttt{low}; classification \texttt{PARTIAL-PROGRESS}, registry status \texttt{partially\_solved}. Created and updated 14 August 2026. These are supplied assessments, not acceptance of a canonical Thurston extension.
\item Attribution: Ulam AI / UnsolvedMath Contributors, CC BY 4.0. The machine-generated research aid is not peer reviewed and requires independent verification before citation.
\end{itemize}
\end{document}
